\honest{Do Scientists Already Know This Stuff?}{Eliezer Yudkowsky}{2008}

A commenter, poke, says that forming good hypotheses is a skill scientists learn in their apprenticeship, and that my description of science is at fault, not science. I did not go through that apprenticeship. But I see plenty of modern scientists making my young self's mistakes, and ``I cannot detect any sign that they were better warned than myself.''

Roger Penrose still thinks consciousness is caused by quantum gravity. ``I expect'' no one warned him against mysterious answers. His view is testable, but it leaves consciousness as mysterious as before. I doubt his thesis advisor warned him, or that anyone warned Niels Bohr. \nb{Penrose came to the view through an argument from Gödel's theorem that human mathematical understanding is not algorithmic. The post does not mention it.} ``As far as I can tell,'' no standard warning exists. I list my four danger signs of a mysterious answer, and say no one ever taught them to me.

The physicists who reject many-worlds were ``Clearly'' never told the formal definition of Occam's razor. \nb{The notes on ``Many Worlds, One Best Guess'' assess that case in full.} Great scientists do tend to have great students, perhaps because mentors pass on skills they cannot describe. But if they cannot put their guidance into words and publish it, ``that's not a good sign for how well these things are understood.'' \nb{Michael Polanyi, whose view is close to poke's, held that the art of research ``cannot be transmitted by prescription, since no prescription for it exists,'' only by example from master to apprentice.}

I doubt that scientists who describe life in 2050 were taught the conjunction fallacy, and whole subfields grow around ``emergence'' and ``complexity.'' Traditional rationality teaches you to give up a theory under crushing evidence. Professional training adds ``Some frequentist stats classes'' and standard techniques. If Science asked more of the average scientist, I do not think Science could get done. \nb{Kuhn had argued something close in 1959: narrow training in a tradition is part of why normal science works. Tversky and Kahneman had found in 1971 that professional psychologists held ``erroneous intuitions about the laws of chance.''}

As Nick Tarleton put it, Science's carefully lax social standard, taken as a private standard, ``lets you believe far too much.'' Many scientists believe ridiculous things outside the laboratory. Is there a standard lecture against that? ``No, as far as I can tell.'' \nb{Feynman said much the same in 1974: the discipline of not fooling yourself was in no course Feynman knew of, and ``We just hope you've caught on by osmosis.''}

A famous mentor might share ``rare personal secrets'': work on the important problems in your field; be more careful than the journal editors demand. \nb{Both had been published in similar words: the first by Richard Hamming in ``You and Your Research'' (1986), the second by Feynman in the same 1974 address.} But there is no secret tradition of precise reasoning, since ``Half of all the scientists out there still believe they believe in God!'' \nb{That is roughly the figure in American surveys. The post gives no evidence that scientists who believe in God only believe that they believe.}
